\section{Mechanism selection: WANT\_SUSPEND}
\label{sec:mechanism}

This is the single most important knob in the policy, and the least obvious.

\subsection{What it actually does}

\kn{WANT\_SUSPEND} does not enable or disable suspension. When the startd decides
a job must go, it decides \textbf{which expression the startd consults}:

\begin{lstlisting}
WANT_SUSPEND true   ->  consults SUSPEND;  PREEMPT is ignored
WANT_SUSPEND false  ->  consults PREEMPT;  SUSPEND is ignored
\end{lstlisting}

\begin{keypoint}
They are \textbf{alternatives, not layers}. A configuration containing a
carefully-written \kn{SUSPEND} and a carefully-written \kn{PREEMPT} runs exactly
one of them, and \kn{WANT\_SUSPEND} picks which.
\end{keypoint}

\subsection{Both failure modes, measured}

Both were found by stress test, and \textbf{neither was visible any other way}.
\kn{condor\_config\_val} printed a correct expression throughout,
\kn{condor\_status} showed healthy machines, and no log recorded anything at all.

\begin{table}[H]
\centering
\begin{tabularx}{\textwidth}{L{0.20\textwidth} L{0.24\textwidth} Y}
\toprule
\textbf{Setting} & \textbf{Consequence} & \textbf{Measured} \\
\midrule
\kn{False} (the default) & \kn{SUSPEND} never fires & Load driven to \textbf{15.72} against a 12.0 threshold and held 90\,s: silent \\
Constant \kn{True} & \kn{PREEMPT} never fires & Memory held \textbf{1.5\,GB below} the eviction floor for 60\,s: silent \\
\bottomrule
\end{tabularx}
\caption{The two ways to get this wrong. The second was introduced while fixing the first.}
\end{table}

The default is \kn{False}, which means a policy that defines \kn{SUSPEND} and
never touches \kn{WANT\_SUSPEND} has a suspension rule that can never execute.

\subsection{The resolution}

It must be an \textbf{expression}, inverting on the pressure that suspension
cannot relieve:

\begin{lstlisting}
EVICT_PRESSURE = memory below floor
              OR memory PSI stalling
              OR disk below floor

WANT_SUSPEND = !(EVICT_PRESSURE)     # CPU contention -> suspend
PREEMPT      =   EVICT_PRESSURE      # memory/disk    -> evict
SUSPEND      = NonCondorLoad > TotalCpus * 0.75
\end{lstlisting}

Read as a sentence: \emph{if the pressure is something freezing the job would
relieve, freeze it; otherwise kill it.} CPU contention is relieved by
\kn{SIGSTOP}. Memory pressure is not, because a stopped process keeps its RSS.

\subsection{Observed behaviour}

Suspension and resumption, with the hysteresis gap doing its job:

\begin{lstlisting}[style=term]
23:34:30  load 11.12                                Busy
23:34:40  SUSPEND is TRUE     Busy -> Suspended         (threshold 12.0)
23:36:45  CONTINUE is TRUE    Suspended -> Busy         (resumed at 7.95)
\end{lstlisting}

Eviction, eleven seconds after pressure began, via the full graceful path:

\begin{lstlisting}[style=term]
00:18:16  PREEMPT is TRUE        Busy -> Retiring
00:18:16  WANT_VACATE is TRUE    Claimed/Retiring -> Preempting/Vacating
00:18:16                         -> Owner/Idle -> Unclaimed -> Delete
00:19:09  job re-matched, evicted again while pressure persisted
\end{lstlisting}

The re-match at \kn{00:19:09} is correct behaviour, not a loop bug: the job
returned to the queue, the negotiator found the machine still advertising
capacity, and the policy evicted it again because the pressure had not lifted.

\subsection{How to test this safely}

\begin{trap}
Do not drive a machine to genuine memory pressure to test a threshold ---
especially not one running a database. The OOM killer does not choose politely,
and there is no equivalent of \kn{renice} for memory.
\end{trap}

\textbf{Raise the floor to meet the memory instead.} Override the eviction
threshold to just under current availability, allocate a couple of gigabytes, and
the real code path runs with minimal risk.

Two details make the difference between a valid test and a misleading one:

\textbf{Override both \kn{PREEMPT} and \kn{WANT\_SUSPEND}.} Overriding
\kn{PREEMPT} alone leaves \kn{WANT\_SUSPEND} pointing at the real floor, where it
stays true, so \kn{SUSPEND} is consulted instead and the test fails for entirely
the wrong reason. This was caught before the first eviction test ran.

\textbf{Generators must be self-terminating and must not fork.} Use \kn{timeout}
to own the deadline, so nothing survives a dropped ssh session or a closed tmux
window.

\begin{trap}
An early CPU load generator tested its own deadline by calling \kn{date +\%s}
each iteration. The fork per iteration dominated the loop, and fourteen workers
produced a load of about 7 --- below the threshold. It looked exactly like a
policy failure. The instrument was broken, not the system.
\end{trap}
